\begin{lemma}[Claim 9: a color unused on the tail]
Fix the Section 7 data, a
\noderef{HierarchyReplacementPresentation} \(P\), and a color \(\alpha\)
missing on its last starred stage and at the final vertex \(y_\ell\).
Assume the ambient Situation 7.1 condition
\[
 \noderef{HierarchyPathNumber}
   (P.\mathit{counterexample}.\mathit{hierarchy})=1.
\]
Assume the paper's normalization: when the divider count is positive,
\(\alpha\) is missing on \(T_{n,q}\), and when it is zero, \(\alpha\) is
missing on \(T_n\).  Assume no index \(1\leq i\leq\ell-1\) is already
common-missing with \(y_\ell\), and assume that no exact displayed tail
edge at or above the starred stage has color \(\alpha\).  Also assume every
proper exactly-one-uncolored replacement coloring has a missing color at
every vertex.

Then \noderef{HierarchyEligibleCommonMissingIndices} is nonempty.
More precisely, the construction produces a replacement presentation in
which \(y_{\ell-1}\) and \(y_\ell\) have a common missing color.
\end{lemma}
\begin{proof}
Write
\[
 \mu_n=P.\mathit{coloring},\qquad
 H=P.\mathit{counterexample}.\mathit{hierarchy},\qquad
 A=\noderef{HierarchyStarredStage}
      (H,\operatorname{Fin.last}(H.\mathit{dividerCount})).
\]
The fields of \noderef{HierarchyReplacementPresentation} give an actual
\noderef{CorrespondingExtendedTashkinovTree}, a hierarchy counterexample,
\noderef{SameGoodHierarchy} data relative to \(C\), and a transported
\noderef{HierarchyMinimumCounterexampleContext}.  In particular,
\(\mu_n\) is \noderef{IsProperEdgeColoring} and
\noderef{IsOneEdgeUncolored} leaves exactly \(e_0\) uncolored.

Unfolding the same-good data gives
\noderef{SameHierarchyStages}.  Its equal divider counts, equal base cores,
and transported equal stage families identify the two last
\noderef{HierarchyStarredStage}s exactly.  Hence the replacement and
original full-tail lengths and prefixes agree.  Put
\[
 \ell=\noderef{HierarchyFullTailLength}(C),\qquad
 B=\noderef{HierarchyFullTailPrefix}(C,\ell-1),
 \qquad y^-=B.\mathit{vertex}(\operatorname{Fin.last}|B|),
 \qquad y=T.\mathit{vertex}(\operatorname{Fin.last}|T|).
\]
Applying \noderef{HierarchyCounterexampleTailAtLeastTwo} and
\noderef{HierarchyCounterexamplePenultimateElementary} to the transported
minimum context yields
\[
 2\leq\ell,\qquad
 \noderef{ColoredVertexSetElementary}(G,\mu_n,V(B)).           \tag{1}
\]
Since \(A\preccurlyeq T\), arithmetic in
\(\ell=|T|-|A|\), together with \(2\leq\ell\), gives
\(|A|+\ell=|T|\).  The defining
\noderef{TreeSequencePrefix} and
\noderef{TreeSequenceInitialSegment} equations therefore give
\[
 |B|=|T|-1,\qquad
 y^-=T.\mathit{vertex}(|T|-1),\qquad A\preccurlyeq B\prec T.  \tag{2}
\]
In particular \(y\notin V(B)\), by vertex injectivity in
\noderef{TreeSequence}.

Apply the stated palette-surplus hypothesis to \(\mu_n\), using the
properness and one-uncolored fields above, and choose
\[
                         \beta\in\overline\mu_n(y^-).          \tag{3}
\]
The integer \(\ell-1\) lies in the interval
\(\{1,\ldots,\ell-1\}\) by (1).  The no-common hypothesis at this exact
index, (3), and
\(\alpha\in\overline\mu_n(y)\) imply
\(\beta\notin\overline\mu_n(y)\), and hence
\[
                              \alpha\ne\beta.                 \tag{4}
\]

We now discharge every first-position premise of
\noderef{ChangeLemma}.  Put
\[
 i_\alpha=\noderef{HierarchyFirstMissingPosition}(\mu_n,T,\alpha),
 \quad i_\beta=\noderef{HierarchyFirstMissingPosition}(\mu_n,T,\beta),
 \quad v_\alpha=T.\mathit{vertex}(i_\alpha),
 \quad v_\beta=T.\mathit{vertex}(i_\beta).
\]
The first part of the start hypothesis gives an \(\alpha\)-missing vertex
in \(A\).  Since \(A\preccurlyeq B\), the least-position definition gives
\[
                 i_\alpha.\mathit{val}\leq |A|.               \tag{5}
\]
By (3), \(\beta\) is missing at the last vertex \(y^-\) of \(B\).
If it were missing at any earlier displayed vertex of \(B\), that vertex
would be distinct from \(y^-\) by \noderef{TreeSequence}, contradicting
the elementarity in (1).  Thus the least-position definition gives
\[
 i_\beta.\mathit{val}=|B|=|T|-1,\qquad v_\beta=y^-.          \tag{6}
\]
The same argument shows that \(\beta\) is missing nowhere on \(A\).
Equations (1), (2), (5), and (6) give
\(i_\alpha.\mathit{val}<i_\beta.\mathit{val}\), so
\(v_\alpha\) precedes \(v_\beta\) in the exact order of \(T\).
Moreover, whenever
\(|A|\leq t<i_\beta.\mathit{val}\), the unused-tail hypothesis gives
\(\mu_n(T.\mathit{edge}(t))\ne\alpha\).  Thus the conditional
order-and-unused-interval premise of \noderef{ChangeLemma} holds: its
antecedent is true because \(\beta\) is not missing on \(A\), its order
conclusion is (5)--(6), and its edge conclusion is the supplied universal
unused-tail assertion.  Both colors are missing on \(B\), by the
\(\alpha\)-witness in \(A\preccurlyeq B\) and (3).  Finally, the
positive-divider or zero-divider case premise is exactly the stated
normalization.  The corresponding conclusion of
\noderef{ChangeLemma} is therefore
\[
 K_\alpha:=\noderef{TwoColorComponent}
   (G,\mu_n,\alpha,\beta,\operatorname{inl}v_\alpha)
 =\noderef{TwoColorComponent}
   (G,\mu_n,\alpha,\beta,\operatorname{inl}v_\beta).          \tag{7}
\]

Let
\[
 Q=\noderef{TwoColorComponent}
      (G,\mu_n,\alpha,\beta,\operatorname{inl}y).
\]
We prove \(Q\ne K_\alpha\), rather than assuming that no component can
meet \(B\).  The vertices \(v_\alpha,v_\beta,y\) are pairwise distinct:
the first two have strictly ordered positions by (5)--(6), and the first
two lie in \(B\) whereas \(y\notin V(B)\) by (2).
By \noderef{MissingColorIncidenceEndpoint}, the vertex-nodes at
\(v_\alpha\) and \(y\) have incidence degree at most one because
\(\alpha\) is missing there.  Applying the same lemma with the two colors
reversed at \(v_\beta\), and using the literal color symmetry of
\noderef{TwoColorIncidenceGraph}, gives the same degree bound there.
If \(Q=K_\alpha\), then (7) puts these three distinct degree-at-most-one
nodes into \(Q\).  But \noderef{TwoColorComponentSpanningPath} realizes
the support of \(Q\) as one simple path, while
\noderef{TwoColorIncidenceDegreeAtMostTwo} bounds every full incidence
degree by two.  Every non-end node of that spanning path already has its
two path neighbors, so a full-degree-at-most-one node must be one of the
path's two ends.  Three distinct such nodes are impossible.  Hence
\(Q\ne K_\alpha\), and the equal-or-disjoint theorem
\noderef{TwoColorComponentsDisjoint} yields
\[
                             Q\cap K_\alpha=\varnothing.       \tag{8}
\]
The missing endpoint \(y\), (4), and
\noderef{MissingColorIncidenceEndpoint} also give a node of degree at most
one in \(Q\); by the definition of
\noderef{TwoColorPathComponent}, \(Q\) is a possibly singleton path
component.

Define the Kempe change
\[
 \sigma_n=\noderef{TwoColorComponentSwap}
   (G,\mu_n,\alpha,\beta,\operatorname{inl}y).                \tag{9}
\]
We first invoke \noderef{StableChangeCasesOneTwo}, before attempting to
fix the exceptional seed.  Supply (1), (4), the two missing witnesses on
\(B\), the path certificate above, (8), the given normalization, and all
properness, one-uncolored, lower-rung, fewer-divider, minimality, and
large-set fields of \(P.\mathit{minimumContext}\).  It gives, in
particular,
\[
       \noderef{HierarchyWeaklyStable}(H,q,\sigma_n),          \tag{10}
\]
where \(q=\operatorname{Fin.last}(H.\mathit{dividerCount})\).
Independently, \noderef{TwoColorComponentSwapProper} and (4) show that
\(\sigma_n\) is proper; (9) fixes uncolored values, so it still leaves
exactly \(e_0\) uncolored.

We next prove that every exact displayed tail-edge value is fixed.  Let
\(t\) satisfy \(|A|\leq t.\mathit{val}<|T|\).  The old value cannot be
\(\alpha\), by the unused-tail hypothesis.  Suppose first that
\(t.\mathit{val}<|T|-1\) and that its old value is \(\beta\).  Identify
this edge with its copy in the prefix \(B\) by (2), and apply
\noderef{HierarchyTailMissingColorExclusion} with \(i=\ell-1\) and the
missing witness (3).  Its conclusion says exactly that the divider count
is zero, \(A=T_n\), \(t.\mathit{val}=|A|\), and this edge is the individual
SE/RE seed \(f_n\), with
\[
 \mu_n(f_n)=\beta\in
 \noderef{TashkinovStoredColors}(P.\mathit{corresponding}.
             \mathit{replacementETT}.\mathit{series}.\mathit{data},n).
                                                                    \tag{11}
\]
It also supplies the actual terminal
\noderef{TashkinovSeriesStep}, its SE-or-RE kind, and the singleton
connecting-edge equation.  Eliminate that certified step.  In the SE
constructor, the singleton equation identifies \(f_n\) with the edge of
the recorded \noderef{MaximumDefectiveChoice}; its vertex-defective and
edge-incidence fields put one end of that exact identity in \(V(T_n)\).
In the RE constructor, it identifies \(f_n\) with the edge of the
\noderef{RevisitingExtensionCertificate}; its inside-current and exact
endpoint fields again put one end in \(V(T_n)\).  Thus in either case
\(f_n\) is a \noderef{TreeIncidentEdge} of \(T_n\).

Because the divider count is zero and the terminal kind is SE or RE,
unfolding \noderef{HierarchyWeaklyStable} in (10) gives the actual
\noderef{StableColoring}
\[
       \sigma_n\text{ is }(T_n,D_n,\mu_n)\text{-stable}.       \tag{12}
\]
The incident-edge clause of (12), applied to (11), now gives
\(\sigma_n(f_n)=\mu_n(f_n)\).  This use of weak stability is prior to and
independent of the full-tail theorem, so there is no circular seed
argument.

It remains to treat the last displayed edge, whose natural position is
\(|T|-1=|A|+\ell-1\).  Since \(\ell\geq2\), the assumed
\noderef{HierarchyPathNumber} equality cannot be its empty-candidate
fallback: that fallback has value \(\ell\ne1\).  Hence the candidate set
is nonempty, its minimum is one, and membership of that minimum in the
defining filter says that every position at least \(|A|+1\) joins its two
consecutive displayed vertices.  Applying this to the last position gives
the exact endpoint equation
\[
 G.\mathit{ends}(T.\mathit{edge}(|T|-1))
       =\{T.\mathit{vertex}(|T|-1),T.\mathit{vertex}(|T|)\}
       =\{y^-,y\}.                                             \tag{13}
\]
By (3) and the definition of \noderef{MissingColors}, no edge incident
with \(y^-\) has color \(\beta\), so (13) excludes \(\beta\) on this last
edge.  The unused-tail hypothesis excludes \(\alpha\).  Thus
\noderef{TwoColorComponentSwap} fixes its value.  Together with the
preceding paragraph and the fact that (9) fixes every value outside
\(\{\alpha,\beta\}\), this proves
\[
 \sigma_n(T.\mathit{edge}(t))=\mu_n(T.\mathit{edge}(t))
       \quad\text{for every }|A|\leq t.\mathit{val}<|T|.       \tag{14}
\]

We now prove, vertex by vertex, the missing-set equality on \(B\):
\[
 \overline{\sigma_n}(v)=\overline{\mu_n}(v)
                         \quad(v\in V(B)).                    \tag{15}
\]
Fix \(v\in V(B)\).  If \(\operatorname{inl}v\in Q\), then (7)--(8)
imply \(v\ne v_\alpha,v_\beta\).  The elementarity in (1), applied to
the missing witnesses \(v_\alpha,v_\beta\in V(B)\), says that neither
\(\alpha\) nor \(\beta\) is missing at \(v\).  Expanding
\noderef{MissingColors} and \noderef{PresentColors}, there is an incident
edge of each affected color.  Properness in the local form of
\noderef{IsProperEdgeColoring} makes each identity unique.  Each of these
two edge-nodes is adjacent to \(\operatorname{inl}v\) in
\noderef{TwoColorIncidenceGraph}, and therefore lies in the same component
\(Q\).  The swap transposes their two values, so both colors remain present
at \(v\); all other colors are fixed.  Hence the missing set at \(v\) is
unchanged.

If \(\operatorname{inl}v\notin Q\), an edge-node of \(Q\) incident with
\(v\) would be adjacent to \(\operatorname{inl}v\) in the incidence graph,
forcing that vertex-node into \(Q\).  Hence no changed edge is incident
with \(v\), and its missing set is again unchanged.  These two cases prove
(15).  In particular, this argument permits internal degree-two vertices
of \(Q\) to lie in \(B\); it does not claim that \(Q\) is vertex-disjoint
from the prefix.

Finally fix
\(\eta\in\noderef{HierarchyStoredColorsAt}(H,q)\).  By the definition of
that stored set, \(\eta\) is not missing on \(A\).  Since \(\alpha\) is
missing on \(A\), \(\eta\ne\alpha\).  If \(\eta\ne\beta\), (9) changes
no edge of color \(\eta\) in either direction, so its missingness at every
vertex, and hence its
\noderef{HierarchyFirstMissingPosition}, is unchanged.  If
\(\eta=\beta\), equation (6) says that \(y^-\) was its first missing
vertex, while (15) says both that it remains missing there and that it is
not missing at any earlier displayed vertex.  Thus its first missing
position is again the position of \(y^-\).  We have proved
\[
 \noderef{HierarchyFirstMissingPosition}(\sigma_n,T,\eta)
 =\noderef{HierarchyFirstMissingPosition}(\mu_n,T,\eta)
 \quad(\eta\in D_{n,q}).                                      \tag{16}
\]
No equality at the terminal vertex was used.

All hypotheses of \noderef{StableChangeCasesOneTwoFullTail} are now
available: (1) and the two prefix witnesses, the path certificate for
\(Q\), (8), the normalized Cases 1/2 alternative, the defining swap (9),
the exact tail equality (14), the vertexwise prefix equality (15), and the
stored first-position equality (16).  Its final conclusion supplies an
actual corresponding ETT on the same exact full tree \(T\) under
\(\sigma_n\), maximum property, a replacement base, a good hierarchy, and
same-good-hierarchy data relative to \(P.\mathit{counterexample}\).

The swap makes \(\beta\) missing at \(y\): originally \(\alpha\) is
missing there and \(\beta\) is present there, and properness makes the
incident \(\beta\)-edge unique; its edge-node is adjacent to the root of
\(Q\), so (9) changes it to \(\alpha\) and leaves no incident
\(\beta\)-edge.  Equation (15) preserves the missingness of \(\beta\) at
\(y^-\).  The two vertices are distinct by (2), so the new full tree is
not elementary.  Together with positivity, maximum property, and the new
good hierarchy, this constructs a genuine
\noderef{HierarchyCounterexample} \(C_\sigma\).

We spell out the transported minimum context.  Rung agreement in the new
\noderef{CorrespondingExtendedTashkinovTree} rewrites the strict
lower-rung bounds, so the lower-rung field of
\(P.\mathit{minimumContext}\) applies verbatim.  Transitivity of the
equalities in \noderef{SameGoodHierarchy} and
\noderef{SameHierarchyStages} gives the same divider count, stages, and
divider positions; hence the fewer-divider field transports and the two
\noderef{HierarchyCounterexampleMeasure}s are equal.  If a same-rung,
equal-divider counterexample were smaller than \(C_\sigma\), it would
therefore be smaller than \(P.\mathit{counterexample}\), contradicting
the old lexicographic-minimality field.  The corresponding certificate's
properness and exact-one-uncolored fields give the first context field.
Its exact terminal ladder-tree agreement and terminal
\noderef{StableColoring} identify the two ladder-tree missing-color sets,
so the old \(2n+11\) inequality gives the new large-set field.  These are
exactly the five fields of
\noderef{HierarchyMinimumCounterexampleContext}.

Compose the new same-good data with \(P.\mathit{sameGood}\), field by
field in \noderef{SameGoodHierarchy}; equality transitivity gives
same-good data relative to the original \(C\).  The new corresponding
ETT, maximum property, base, \(C_\sigma\), composed same-good witness, and
transported minimum context are precisely a new
\noderef{HierarchyReplacementPresentation} \(P'\).

Take \(i=\ell-1\), \(P'\), and \(\beta\) in the classical filter defining
\noderef{HierarchyEligibleCommonMissingIndices}.  The interval bounds
follow from \(2\leq\ell\); equation (15) at the last vertex of \(B\) and
the preceding endpoint calculation give the two required missing-color
memberships.  Therefore the exact index \(\ell-1\) belongs to that finite
set, as claimed.
\end{proof}
